\clearpage
\SourceLeaf{018}\pdfbookmark[0]{Contenta.}{contenta}
\begin{center}
{\Large CONTENTA.}\par\bigskip
{\large LIBER PRIMUS.}\par\medskip
Relationes generales inter quantitates, per quas corporum coelestium motus circa Solem definiuntur.
\end{center}
\noindent\begin{tabular}{@{}p{.18\linewidth}p{.78\linewidth}@{}}
Sectio I. & \emph{Relationes ad locum simplicem in orbita spectantes.}\\[.4em]
Sectio II. & \emph{Relationes ad locum simplicem in spatio spectantes.}\\[.4em]
Sectio III. & \emph{Relationes inter locos plures in orbita.}\\[.4em]
Sectio IV. & \emph{Relationes inter locos plures in spatio.}
\end{tabular}
\begin{center}\bigskip
{\large LIBER SECUNDUS.}\par\medskip
Investigatio orbitarum corporum coelestium ex observationibus geocentricis.
\end{center}
\noindent\begin{tabular}{@{}p{.18\linewidth}p{.78\linewidth}@{}}
Sectio I. & \emph{Determinatio orbitae e tribus observationibus completis.}\\[.4em]
Sectio II. & \emph{Determinatio orbitae e quatuor observationibus, quarum duae tantum completae sunt.}\\[.4em]
Sectio III. & \emph{Determinatio orbitae observationibus quotcunque quam proxime satisfacientis.}\\[.4em]
Sectio IV. & \emph{De determinatione orbitarum, habita ratione perturbationum.}
\end{tabular}
\par\bigskip\noindent\textsc{Tabulae.}
